\documentclass[10pt]{article}

% ========== Packages ==========
\usepackage[utf8]{inputenc}
\usepackage{amsmath, amssymb}
\usepackage{graphicx}
\usepackage{booktabs}
\usepackage{array}
\usepackage{hyperref}
\usepackage{geometry}
\geometry{margin=1in}

% ========== Title ==========
\title{DVCL: Dual-View Contrastive Learning for Robust Heterogeneous Graph Neural Networks}

\author{Anonymous Authors}

\date{}

\begin{document}

\maketitle

% =========================
% Abstract
% =========================
\begin{abstract}
Heterogeneous graph neural networks have achieved promising performance in modeling complex relational data with multiple types of nodes and edges.
However, they remain vulnerable to structural adversarial attacks, where a small number of malicious edge perturbations may be amplified through meta-path-based message propagation and further corrupt high-order semantic neighborhoods.
Existing robust heterogeneous graph methods mainly defend against such perturbations within the topological semantic space, for example by estimating edge confidence, filtering suspicious meta-path-induced edges, or learning semantic attention weights.
Although effective to some extent, these methods still heavily rely on the attacked heterogeneous topology to construct candidate neighborhoods, making it difficult to recover reliable representations when the topology itself has been severely polluted.

To address this limitation, we propose DVCL, a dual-view contrastive learning framework for robust heterogeneous graph learning.
DVCL constructs two complementary views: a purified topology view derived from meta-path-induced semantic structures, and a feature-induced view built from node feature similarity.
The topology view preserves heterogeneous structural semantics through meta-path modeling, semantic attention, and two-stage edge filtering, while the feature view provides a relatively topology-independent reference under structural attacks.
To further reduce semantic inconsistency between the two views, DVCL employs a symmetric cross-view contrastive objective that aligns the topology representation and feature representation of the same node while distinguishing representations from different nodes.
Together with supervised node classification and an auxiliary semantic topology learning objective, DVCL learns robust node representations that integrate heterogeneous semantic expressiveness and feature-level stability.
Experiments on benchmark heterogeneous graphs under structural attacks demonstrate that DVCL improves node classification robustness compared with existing heterogeneous graph defense baselines.
\end{abstract}

% =========================
% 1 Introduction
% =========================
\section{Introduction}

Heterogeneous graphs are widely used to model real-world systems that contain multiple types of entities and relations, such as academic networks, recommendation platforms, biomedical networks, and social media systems. Compared with homogeneous graphs, heterogeneous graphs provide richer semantic information through relation types and meta-paths. This semantic structure enables heterogeneous graph neural networks (HGNNs), such as HAN, HGT, MAGNN, and HeCo, to learn expressive node representations by aggregating information from multiple semantic neighborhoods \cite{wang2019han,hu2020hgt,fu2020magnn,wang2021heco}. These models have achieved strong performance on node classification and representation learning tasks.

Despite their effectiveness, HGNNs remain vulnerable to adversarial structural perturbations. A small number of malicious edge insertions or deletions can change not only local neighborhoods, but also the high-order semantic neighborhoods induced by meta-paths. This issue is particularly serious in heterogeneous graphs because a perturbation on one relation type may be amplified after meta-path composition and may affect many target-type nodes. As a result, semantic aggregation may propagate misleading information, and attention mechanisms may assign large weights to corrupted semantic views.

Existing robust graph learning methods mitigate structural attacks through edge pruning, probabilistic modeling, graph structure learning, or similarity-based message passing \cite{wu2019gcnjaccard,zhu2019rgcn,jin2020prognn,zhang2020gnnguard,jin2021simpgcn}. Recent robust heterogeneous graph methods further consider typed relations and semantic structures, for example by estimating edge confidence or suppressing suspicious meta-path-induced connections \cite{rohe,hseco,heteguard}. However, most of these methods still construct candidate neighborhoods from the observed heterogeneous topology. When the topology is heavily polluted, topology-only defenses may have limited ability to recover reliable representations because their denoising process is still constrained by attacked structural evidence.

To address this limitation, we propose DVCL, a Dual-View Contrastive Learning framework for robust heterogeneous graph neural networks. DVCL constructs a purified topology view from meta-path-induced semantic graphs and a feature-induced view from target-node feature similarity. The topology view preserves heterogeneous semantic expressiveness through meta-path modeling, semantic attention, and two-stage edge filtering. The feature-induced view provides a complementary, relatively topology-independent reference under structural attacks. DVCL further aligns the two views with a symmetric cross-view contrastive objective, encouraging each node to maintain consistent representations across semantic topology and feature similarity while remaining discriminative from other nodes.

The main contributions of this paper are summarized as follows:
\begin{itemize}
    \item We propose DVCL, a dual-view robust heterogeneous graph learning framework that jointly exploits purified semantic topology and feature-induced similarity for node classification under structural attacks.
    \item We design a topology purification strategy that combines meta-path-level edge confidence, semantic attention, and global semantic filtering, and pair it with a feature-induced KNN view to reduce reliance on attacked graph structure.
    \item We introduce a symmetric cross-view contrastive objective and a complete evaluation protocol for clean, global-attack, and target-attack settings on heterogeneous graph benchmarks.
\end{itemize}

\section{Preliminaries}
\label{sec:preliminaries}

\noindent\textbf{Definition 1. Heterogeneous Graph.}
A heterogeneous graph is defined as
\begin{equation}
\mathcal{G}=(\mathcal{V},\mathcal{E},\mathbf{X},\phi,\psi),
\end{equation}
where $\mathcal{V}$ and $\mathcal{E}$ denote the node set and edge set, respectively.
$\mathbf{X}\in\mathbb{R}^{N\times d}$ is the node feature matrix.
The mapping function $\phi:\mathcal{V}\rightarrow\mathcal{A}$ assigns each node to a node type, and $\psi:\mathcal{E}\rightarrow\mathcal{R}$ assigns each edge to a relation type.
A graph is heterogeneous when it contains multiple types of nodes or edges.

\noindent\textbf{Definition 2. Meta-path.}
A meta-path is a composite relation defined on the schema of a heterogeneous graph:
\begin{equation}
\mathcal{P}:
\mathcal{T}_1 \xrightarrow{\mathcal{R}_1}
\mathcal{T}_2 \xrightarrow{\mathcal{R}_2}
\cdots
\xrightarrow{\mathcal{R}_L}
\mathcal{T}_{L+1},
\end{equation}
where $\mathcal{T}_l\in\mathcal{A}$ denotes a node type and $\mathcal{R}_l\in\mathcal{R}$ denotes a relation type.
A meta-path describes a specific semantic relation between nodes.
For example, in an academic network, author--paper--author connects authors who write the same paper, while author--paper--venue--paper--author connects authors through publication venues.

\noindent\textbf{Definition 3. Meta-path-based Neighbors.}
Given a meta-path $\mathcal{P}$, the meta-path-based neighbor set of node $v_i$ is defined as
\begin{equation}
\mathcal{N}_{\mathcal{P}}(v_i)
= \left\{ v_j \in \mathcal{V} \mid v_i \xrightarrow{\mathcal{P}} v_j \right\}, 
\end{equation}
where $v_i \xrightarrow{\mathcal{P}} v_j$ means that node $v_i$ can reach node $v_j$ through at least one path instance following $\mathcal{P}$.
Different meta-paths produce different semantic neighborhoods for the same node.

\noindent\textbf{Definition 4. Meta-path-based Semantic Aggregation.}
Meta-path-based heterogeneous graph learning usually first constructs multiple semantic views according to different meta-paths.
For each meta-path, a node aggregates information from its corresponding meta-path-based neighbors.
Then, representations from different meta-paths are combined by semantic-level aggregation.
This process can be generally described as
\begin{equation}
\mathbf{h}_i
=
\operatorname{Fuse}
\left(
\mathbf{h}_i^{(1)},\mathbf{h}_i^{(2)},\ldots,\mathbf{h}_i^{(M)}
\right),
\end{equation}
where $\mathbf{h}_i^{(m)}$ denotes the representation of node $v_i$ under the $m$-th meta-path, and $\operatorname{Fuse}(\cdot)$ denotes a semantic fusion function.
In this paper, such semantic aggregation provides the basis for constructing the topology view.

\noindent\textbf{Definition 5. Adversarial Structural Attack.}
Given a clean heterogeneous graph $\mathcal{G}$, an adversarial structural attack generates a perturbed graph
\begin{equation}
\widetilde{\mathcal{G}}
=
(\mathcal{V},\widetilde{\mathcal{E}},\mathbf{X},\phi,\psi),
\end{equation}
by modifying the edge set under a perturbation budget:
\begin{equation}
|\widetilde{\mathcal{E}}\triangle\mathcal{E}| \leq \Delta,
\end{equation}
where $\triangle$ denotes the symmetric difference between two edge sets and $\Delta$ is the attack budget.
A global attack aims to degrade the overall classification performance of the model, while a targeted attack aims to mislead the predictions of specific target nodes.
In heterogeneous graphs, structural perturbations may further affect high-order semantic neighborhoods induced by meta-paths.



% ========== 3 Method ==========
\section{Method}

\subsection{Overview}

We propose Dual-View Contrastive Learning (DVCL) for robust heterogeneous graph node classification. Given a heterogeneous graph
$\mathcal{G}=(\mathcal{V},\mathcal{E},\mathcal{A},\mathcal{R})$, where $\mathcal{A}$ and $\mathcal{R}$ denote the node-type and relation-type sets, respectively, our goal is to predict the labels of target-type nodes $\mathcal{V}_t=\{v_1,\ldots,v_n\}$. Let $X\in\mathbb{R}^{n\times d}$ be the feature matrix of the target nodes and $Y$ be the label set.

DVCL constructs two complementary views for the target nodes. The first one is a purified topology view, which is obtained from meta-path-induced homogeneous graphs and semantic attention. The second one is a feature-induced view, which is built by a KNN graph in an enhanced feature space. The two views are encoded by separate graph encoders, fused for node classification, and further aligned by a symmetric cross-view InfoNCE loss.

\subsection{Purified Topology View Construction}

\paragraph{Meta-path-induced graph construction.}
Let $\mathcal{P}=\{P_1,\ldots,P_M\}$ denote the set of meta-paths, where $M$ is the number of meta-paths. For each meta-path $P_m$, we construct a target-node homogeneous adjacency matrix $A^{(m)}\in\mathbb{R}^{n\times n}$ by multiplying the row-normalized relation-specific adjacency matrices along the meta-path:
\begin{equation}
A^{(m)}
=
\bar{A}_{r_1}\bar{A}_{r_2}\cdots \bar{A}_{r_L},
\end{equation}
where $P_m=(r_1,r_2,\ldots,r_L)$ and $\bar{A}_{r_l}$ is the row-normalized adjacency matrix of relation $r_l$.

\paragraph{Edge confidence estimation.}
The raw meta-path connectivity only reflects structural reachability. To further incorporate node feature consistency, we compute a feature-calibrated edge score for each candidate edge $(i,j)$ under meta-path $P_m$:
\begin{equation}
s_{ij}^{(m)}
=
A_{ij}^{(m)} \cdot x_i^{\top}x_j,
\end{equation}
where $x_i$ and $x_j$ are the features of target nodes $v_i$ and $v_j$. We then apply degree normalization:
\begin{equation}
w_{ij}^{(m)}
=
(d_i^{(m)})^{-1}
s_{ij}^{(m)}
(d_j^{(m)})^{-1},
\end{equation}
where $d_i^{(m)}=\sum_j s_{ij}^{(m)}$. The value $w_{ij}^{(m)}$ is used as the confidence score of edge $(i,j)$ in the $m$-th meta-path graph.

\paragraph{Meta-path-level filtering.}
For each meta-path graph, we remove low-confidence edges by a meta-path-specific threshold $\gamma_m$:
\begin{equation}
E^{(m)}_{\mathrm{pur}}
=
\{(i,j)\mid w_{ij}^{(m)} \geq \gamma_m\}.
\end{equation}
This produces a set of purified meta-path graphs
$\{\mathcal{G}^{(1)}_{\mathrm{pur}},\ldots,\mathcal{G}^{(M)}_{\mathrm{pur}}\}$.

\paragraph{Semantic attention.}
We use a HAN-style semantic attention module to estimate the importance of different meta-paths. For each purified meta-path graph, a GAT encoder produces a meta-path-specific node representation $h_i^{(m)}$. The semantic score of meta-path $P_m$ is computed by
\begin{equation}
e_m
=
\frac{1}{n}
\sum_{i=1}^{n}
\mathbf{q}^{\top}
\tanh(W_s h_i^{(m)} + b_s),
\end{equation}
where $W_s$, $b_s$, and $\mathbf{q}$ are learnable parameters. The normalized semantic attention weight is
\begin{equation}
\beta_m
=
\frac{\exp(e_m)}
{\sum_{m'=1}^{M}\exp(e_{m'})}.
\end{equation}

\paragraph{Fused topology graph purification.}
The final topology view is constructed by aggregating the purified meta-path edge scores with semantic attention:
\begin{equation}
\bar{w}_{ij}
=
\sum_{m=1}^{M}
\beta_m w_{ij}^{(m)}.
\end{equation}
After coalescing duplicated edges from different meta-paths, we apply a global semantic filtering threshold $\gamma_s$:
\begin{equation}
E_{\mathrm{topo}}
=
\{(i,j)\mid \bar{w}_{ij}\geq \gamma_s\}.
\end{equation}
The resulting graph is denoted as
$\mathcal{G}_{\mathrm{topo}}=(\mathcal{V}_t,E_{\mathrm{topo}})$ and serves as the purified topology view.

\subsection{Feature-Induced View Construction}

The topology view is constructed from meta-path-induced structures and may still be affected by adversarial perturbations on graph edges. To provide a complementary view that is less dependent on the observed topology, DVCL constructs a feature-induced graph based only on the raw attributes of target nodes.

Let $X=[x_1,x_2,\ldots,x_n]^\top\in\mathbb{R}^{n\times d}$ denote the raw feature matrix of the target nodes, where $x_i$ is the feature vector of node $v_i$. We first apply $\ell_2$ normalization to each feature vector:
\begin{equation}
\tilde{x}_i
=
\frac{x_i}{\|x_i\|_2}.
\end{equation}
Then we compute the cosine similarity between each pair of target nodes:
\begin{equation}
\operatorname{sim}(i,j)
=
\tilde{x}_i^\top \tilde{x}_j.
\end{equation}
For each node $v_i$, we select the top-$k$ most similar nodes as its feature neighbors:
\begin{equation}
\mathcal{N}_k(i)
=
\operatorname{TopK}_{j\neq i}
\operatorname{sim}(i,j).
\end{equation}
The feature-induced edge set is defined as
\begin{equation}
E_{\mathrm{feat}}
=
\{(i,j)\mid j\in\mathcal{N}_k(i)\}.
\end{equation}
This produces a directed KNN graph
\begin{equation}
\mathcal{G}_{\mathrm{feat}}
=
(\mathcal{V}_t,E_{\mathrm{feat}}),
\end{equation}
which serves as the feature-induced view. Since this view is built from raw node attributes rather than the perturbed heterogeneous topology, it provides a complementary signal to the purified topology view under structural attacks.

\subsection{Dual-View Encoding}

DVCL uses two separate graph encoders for the topology view and the feature-induced view. Both encoders are implemented as GAT-based graph encoders with self-loops.

For the topology view, the node representation is
\begin{equation}
Z_{\mathrm{topo}}
=
f_{\mathrm{topo}}(X,\mathcal{G}_{\mathrm{topo}}),
\end{equation}
where $f_{\mathrm{topo}}(\cdot)$ denotes the topology-view encoder and
$Z_{\mathrm{topo}}\in\mathbb{R}^{n\times d_z}$.

For the feature view, we first apply feature dropout to the input feature matrix and then encode it on the feature-induced graph:
\begin{equation}
Z_{\mathrm{feat}}
=
f_{\mathrm{feat}}(\operatorname{Dropout}(X),\mathcal{G}_{\mathrm{feat}}),
\end{equation}
where $f_{\mathrm{feat}}(\cdot)$ is the feature-view encoder and
$Z_{\mathrm{feat}}\in\mathbb{R}^{n\times d_z}$. The feature masking/dropout operation serves as a lightweight regularizer for the feature view.

\subsection{Dual-View Fusion and Classification}

For each target node $v_i$, DVCL obtains two view-specific embeddings
$z_i^{\mathrm{topo}}$ and $z_i^{\mathrm{feat}}$. In the default setting, the two embeddings are concatenated:
\begin{equation}
z_i
=
z_i^{\mathrm{topo}}
\Vert
z_i^{\mathrm{feat}}.
\end{equation}
The final prediction logits are computed by a linear classifier:
\begin{equation}
o_i
=
W_c z_i + b_c,
\end{equation}
where $W_c$ and $b_c$ are learnable classifier parameters.

Given the labeled training node set $\mathcal{V}_{\mathrm{tr}}$, the classification loss is
\begin{equation}
\mathcal{L}_{\mathrm{cls}}
=
-\frac{1}{|\mathcal{V}_{\mathrm{tr}}|}
\sum_{i\in\mathcal{V}_{\mathrm{tr}}}
\log
\frac{\exp(o_{i,y_i})}
{\sum_{c=1}^{C}\exp(o_{i,c})},
\end{equation}
where $C$ is the number of classes and $y_i$ is the ground-truth label of node $v_i$.

\subsection{Cross-View Contrastive Learning}

To encourage consistent representations between the topology view and the feature view, DVCL applies a symmetric cross-view InfoNCE loss on training nodes. We first normalize the two view-specific embeddings:
\begin{equation}
\tilde{z}_i^{\mathrm{topo}}
=
\frac{z_i^{\mathrm{topo}}}{\|z_i^{\mathrm{topo}}\|_2},
\quad
\tilde{z}_i^{\mathrm{feat}}
=
\frac{z_i^{\mathrm{feat}}}{\|z_i^{\mathrm{feat}}\|_2}.
\end{equation}
For node $v_i$, the positive pair is
$(\tilde{z}_i^{\mathrm{topo}},\tilde{z}_i^{\mathrm{feat}})$, while embeddings of other nodes in the opposite view are treated as negatives. The topology-to-feature contrastive loss is
\begin{equation}
\mathcal{L}_{t\rightarrow f}
=
-\frac{1}{|\mathcal{V}_{\mathrm{tr}}|}
\sum_{i\in\mathcal{V}_{\mathrm{tr}}}
\log
\frac{
\exp((\tilde{z}_i^{\mathrm{topo}})^{\top}
\tilde{z}_i^{\mathrm{feat}}/\tau)
}
{
\sum_{j\in\mathcal{V}_{\mathrm{tr}}}
\exp((\tilde{z}_i^{\mathrm{topo}})^{\top}
\tilde{z}_j^{\mathrm{feat}}/\tau)
},
\end{equation}
where $\tau$ is the temperature parameter. Similarly, the feature-to-topology loss is
\begin{equation}
\mathcal{L}_{f\rightarrow t}
=
-\frac{1}{|\mathcal{V}_{\mathrm{tr}}|}
\sum_{i\in\mathcal{V}_{\mathrm{tr}}}
\log
\frac{
\exp((\tilde{z}_i^{\mathrm{feat}})^{\top}
\tilde{z}_i^{\mathrm{topo}}/\tau)
}
{
\sum_{j\in\mathcal{V}_{\mathrm{tr}}}
\exp((\tilde{z}_i^{\mathrm{feat}})^{\top}
\tilde{z}_j^{\mathrm{topo}}/\tau)
}.
\end{equation}
The final cross-view contrastive loss is
\begin{equation}
\mathcal{L}_{\mathrm{cl}}
=
\frac{1}{2}
(
\mathcal{L}_{t\rightarrow f}
+
\mathcal{L}_{f\rightarrow t}
).
\end{equation}

\subsection{Training Objective}

In addition to the final classification loss and cross-view contrastive loss, DVCL keeps the auxiliary HAN branch classification loss to directly supervise the semantic attention module. Let $o_i^{\mathrm{HAN}}$ be the logits produced by the HAN semantic branch. The auxiliary HAN loss is
\begin{equation}
\mathcal{L}_{\mathrm{HAN}}
=
-\frac{1}{|\mathcal{V}_{\mathrm{tr}}|}
\sum_{i\in\mathcal{V}_{\mathrm{tr}}}
\log
\frac{\exp(o^{\mathrm{HAN}}_{i,y_i})}
{\sum_{c=1}^{C}\exp(o^{\mathrm{HAN}}_{i,c})}.
\end{equation}
The overall training objective is
\begin{equation}
\mathcal{L}
=
\lambda_{\mathrm{HAN}}\mathcal{L}_{\mathrm{HAN}}
+
\mathcal{L}_{\mathrm{cls}}
+
\lambda_{\mathrm{DVCL}}\mathcal{L}_{\mathrm{cl}},
\end{equation}
where $\lambda_{\mathrm{HAN}}$ controls the strength of the auxiliary HAN supervision and $\lambda_{\mathrm{DVCL}}$ controls the strength of cross-view contrastive learning. All trainable parameters, including the semantic attention module, the two view encoders, and the final classifier, are optimized jointly.

% ========== 4 Experiments ==========
% ========== 4 Experiments ==========
\section{Experiments}

\subsection{Datasets}

We evaluate DVCL on four heterogeneous graph benchmarks: ACM, DBLP, AMiner, and IMDB. Each dataset contains multiple node types and relation types, and the evaluation task is semi-supervised node classification on a target node type. Following the HSeCo setting, we use dataset-specific meta-paths to construct semantic neighborhoods for target nodes. Table~\ref{tab:dataset_settings} summarizes the target node type and meta-paths used in each dataset.

\begin{table}[t]
\centering
\caption{Dataset settings. Detailed statistics such as the number of nodes, edges, classes, and train/validation/test splits will be filled after the final data preprocessing is fixed.}
\label{tab:dataset_settings}
\begin{tabular}{c|c|c}
\hline
Dataset & Target node type & Meta-paths \\
\hline
ACM & Paper & P-A-P, P-F-P \\
DBLP & Author & A-P-A, A-P-C-P-A, A-P-T-P-A \\
AMiner & Paper & P-A-P, P-R-P \\
IMDB & Movie & M-D-M, M-A-M \\
\hline
\end{tabular}
\end{table}

ACM and DBLP represent academic networks, where papers, authors, venues, fields, conferences, and terms form typed semantic relations. AMiner is also an academic network but contains different relation patterns and label distributions. IMDB is a movie network in which movies are connected through actors and directors. These datasets allow us to test whether DVCL can improve robustness across different graph schemas and target node types.

\subsection{Baselines}

We compare DVCL with representative heterogeneous graph neural networks, robust graph neural networks, and robust heterogeneous graph learning methods. The heterogeneous graph baselines include HAN \cite{wang2019han}, HGT \cite{hu2020hgt}, MAGNN \cite{fu2020magnn}, HeCo \cite{wang2021heco}, HeteroSAGE, and SeHGNN. These methods evaluate whether DVCL improves robustness over strong semantic aggregation and representation learning models.

The robust graph baselines include GCN-Jaccard \cite{wu2019gcnjaccard}, RGCN \cite{zhu2019rgcn}, Pro-GNN \cite{jin2020prognn}, GNNGuard \cite{zhang2020gnnguard}, and SimP-GCN \cite{jin2021simpgcn}, when they can be adapted to the target-node homogeneous graphs induced by meta-paths. We also compare with robust heterogeneous graph methods, including RoHe \cite{rohe}, HeteGuard \cite{heteguard}, and HSeCo \cite{hseco}. HSeCo is the most direct baseline because DVCL builds on the idea of semantic topology learning while adding an explicit feature-induced view and cross-view alignment.

\subsection{Attack Settings}

We consider both global structural attacks and target structural attacks. For global attacks, we evaluate PRBCD and HetePRBCD with perturbation rates of $5\%$, $10\%$, $15\%$, $20\%$, and $25\%$. PRBCD performs scalable global adversarial structure perturbation, while HetePRBCD extends the attack setting to heterogeneous relations. We also report clean performance without attack.

For target attacks, we use an FGSM-based heterogeneous graph target attack. For each random seed, 500 target nodes are randomly sampled from the test set. Unless otherwise specified, the attack budget is set to 3. The attacker perturbs task-relevant relations for each dataset, such as P-A edges on ACM, so that the perturbation can influence the meta-path-induced semantic neighborhoods of target nodes.

\subsection{Implementation Details}

All models are trained under the same data splits and attack settings for fair comparison. We run each experiment with five random seeds, namely 1, 2, 3, 4, and 5, and report the mean and standard deviation. The maximum number of training epochs is set to 200, and early stopping is applied with patience 100. We use Adam as the optimizer. Unless otherwise specified, the learning rate is set to $0.005$, the weight decay is set to $0.001$, and the dropout rate is set to $0.3$.

DVCL contains a purified topology view and a feature-induced view. The topology view is constructed from purified meta-path graphs with semantic attention. The feature view is constructed by a KNN graph over target-node representations. The two view-specific representations are concatenated for final prediction. The feature masking rate is set to $0.2$, and the KNN size is set to $k=20$. Based on validation performance, we set $\lambda_{\mathrm{DVCL}}=0.5$ and $\lambda_{\mathrm{HAN}}=1.0$. For the optimized feature view, we construct the KNN graph using the concatenation of raw node features and normalized relation-wise neighbor feature means, and retain mutual KNN edges to reduce noisy feature-view connections.

For global attacks, we report Micro-F1 and Macro-F1 on the test set. For single-label multi-class classification datasets, accuracy is equivalent to Micro-F1. For target attacks, we report clean target Micro-F1/Macro-F1, attacked target Micro-F1/Macro-F1, and the performance drop after attack.

\subsection{Overall Performance}

Table~\ref{tab:overall_global} reports the current five-seed ACM results under global structural attacks. This table uses the optimized feature-view construction, denoted as mutual neighbor-mean KNN, where the KNN graph is constructed from raw target-node features augmented with normalized relation-wise neighbor feature means and then filtered by mutual nearest-neighbor consistency. All values are reported as mean $\pm$ standard deviation over seeds 1--5.

\begin{table}[t]
\centering
\caption{ACM global attack performance of DVCL with mutual neighbor-mean feature view.}
\label{tab:overall_global}
\begin{tabular}{c|c|c|c|c}
\hline
Dataset & Attack & Rate & Micro-F1 & Macro-F1 \\
\hline
ACM & Clean & 0\% & $0.8966\pm0.0043$ & $0.8949\pm0.0042$ \\
ACM & PRBCD & 5\% & $0.8913\pm0.0059$ & $0.8899\pm0.0052$ \\
ACM & PRBCD & 15\% & $0.8838\pm0.0016$ & $0.8820\pm0.0028$ \\
ACM & PRBCD & 25\% & $0.8785\pm0.0024$ & $0.8777\pm0.0015$ \\
ACM & HetePRBCD & 5\% & $0.8930\pm0.0018$ & $0.8918\pm0.0026$ \\
ACM & HetePRBCD & 15\% & $0.8920\pm0.0019$ & $0.8902\pm0.0018$ \\
ACM & HetePRBCD & 25\% & $0.8893\pm0.0062$ & $0.8876\pm0.0058$ \\
\hline
\end{tabular}
\end{table}

The results show that DVCL maintains stable performance as the attack budget increases. The degradation from clean to PRBCD 25\% is $1.81$ percentage points in Micro-F1, while the degradation from clean to HetePRBCD 25\% is $0.73$ percentage points. This suggests that the optimized feature-induced view can provide useful complementary information when heterogeneous structural semantics are perturbed.

We also found earlier single-seed DBLP results in the experiment notes. Under the default DVCL v0 setting, DVCL improves DBLP HetePRBCD Micro-F1 from $0.8051$ to $0.8331$ at 5\%, from $0.7813$ to $0.8198$ at 10\%, from $0.7336$ to $0.8094$ at 15\%, from $0.7370$ to $0.8121$ at 20\%, and from $0.7321$ to $0.8174$ at 25\%. Since these DBLP values are single-seed results, we use them as supporting evidence rather than mixing them into the five-seed ACM main table.
\subsection{Ablation Study}

Table~\ref{tab:ablation} reports the standard ACM component ablation under a controlled raw KNN setting. In this ablation, all variants use raw feature KNN with $k=20$ so that the effects of the dual-view structure and the cross-view contrastive loss can be isolated. The attack average is computed over PRBCD and HetePRBCD at 5\%, 15\%, and 25\%.

\begin{table}[t]
\centering
\caption{ACM component ablation under raw KNN feature-view construction.}
\label{tab:ablation}
\begin{tabular}{c|c|c|c}
\hline
Variant & Clean Micro-F1 & Attack Micro-F1 Avg & Attack Macro-F1 Avg \\
\hline
Full DVCL & $0.8770\pm0.0071$ & 0.8710 & 0.8680 \\
DVCL w/o CL & $0.8743\pm0.0074$ & 0.8656 & 0.8626 \\
Topology only & $0.8523\pm0.0177$ & 0.8471 & 0.8415 \\
Feature only & $0.8539\pm0.0086$ & 0.8539 & 0.8512 \\
\hline
\end{tabular}
\end{table}

Full DVCL is the best variant on clean graphs and on the attack average. Removing cross-view contrastive learning reduces the attack-average Micro-F1 by $0.53$ percentage points, showing that InfoNCE contributes beyond simple representation concatenation. The two single-view variants are both weaker than the full model, which confirms that topology-view semantics and feature-view stability are complementary.

We further evaluate feature-view construction strategies in Table~\ref{tab:feature_ablation}. Compared with raw KNN, adding heterogeneous neighbor information improves both clean and attacked performance. The strongest current candidate is mutual neighbor-mean KNN, which improves the attack-average Micro-F1 from 0.8710 to 0.8880.

\begin{table}[t]
\centering
\caption{ACM feature-view construction ablation.}
\label{tab:feature_ablation}
\begin{tabular}{c|c|c|c}
\hline
Feature view & Clean Micro-F1 & Attack Micro-F1 Avg & Attack Macro-F1 Avg \\
\hline
Raw KNN & 0.8770 & 0.8710 & 0.8680 \\
Mutual hetero-profile KNN & 0.8897 & 0.8806 & 0.8787 \\
Neighbor-mean KNN & 0.8933 & 0.8826 & 0.8811 \\
Mutual neighbor-mean KNN & 0.8966 & 0.8880 & 0.8866 \\
\hline
\end{tabular}
\end{table}
\subsection{Parameter Sensitivity}

We analyze the sensitivity of DVCL to the cross-view contrastive loss weight $\lambda_{\mathrm{DVCL}}$. Table~\ref{tab:param} reports five-seed ACM results under raw KNN feature-view construction. The attack average is computed over PRBCD and HetePRBCD at 5\%, 15\%, and 25\%.

\begin{table}[t]
\centering
\caption{Sensitivity to $\lambda_{\mathrm{DVCL}}$ on ACM.}
\label{tab:param}
\begin{tabular}{c|c|c|c}
\hline
$\lambda_{\mathrm{DVCL}}$ & Clean Micro-F1 & Attack Micro-F1 Avg & Attack Macro-F1 Avg \\
\hline
0.5 & $0.8770\pm0.0071$ & 0.8710 & 0.8680 \\
1.0 & $0.8810\pm0.0049$ & 0.8723 & 0.8690 \\
2.0 & $0.8790\pm0.0085$ & 0.8734 & 0.8709 \\
\hline
\end{tabular}
\end{table}

Increasing $\lambda_{\mathrm{DVCL}}$ from 0.5 to 1.0 improves clean Micro-F1 and slightly improves the attack average. A larger value of 2.0 obtains the highest attack-average Micro-F1, but the experiment notes show larger seed variance on HetePRBCD 15\% and 25\%. We therefore treat $\lambda_{\mathrm{DVCL}}=1.0$ as a balanced candidate and $\lambda_{\mathrm{DVCL}}=2.0$ as a stronger but less stable robust setting.
\subsection{Robustness Analysis}

We further evaluate robustness under an FGSM-style target attack on ACM. Each seed samples 500 target paper nodes from the test set, and the attack budget is set to 3. Table~\ref{tab:target_attack} reports five-seed target-node results for different feature-view constructions.

\begin{table}[t]
\centering
\caption{ACM target attack robustness under different feature-view constructions.}
\label{tab:target_attack}
\begin{tabular}{c|c|c|c|c}
\hline
Feature view & Clean Micro-F1 & Attacked Micro-F1 & Drop & Attacked Macro-F1 \\
\hline
Raw KNN & $0.8736\pm0.0120$ & $0.8656\pm0.0184$ & $-0.80\pm0.93$ pp & $0.8632\pm0.0193$ \\
Hetero-profile KNN & $0.8788\pm0.0150$ & $0.8684\pm0.0169$ & $-1.04\pm0.65$ pp & $0.8667\pm0.0197$ \\
Mutual hetero-profile KNN & $0.8852\pm0.0190$ & $0.8784\pm0.0217$ & $-0.68\pm0.58$ pp & $0.8778\pm0.0227$ \\
\hline
\end{tabular}
\end{table}

The mutual hetero-profile feature view obtains the best attacked target Micro-F1 and Macro-F1, improving attacked target Micro-F1 by $1.28$ percentage points over raw KNN. Its target drop is also slightly smaller than raw KNN, indicating that the gain does not only come from higher clean target performance. Together with the global attack results, this supports the design choice of enriching the feature-induced view with heterogeneous neighborhood information.
% ========== 5 Related Work ==========
\section{Related Work}

\subsection{Heterogeneous Graph Neural Networks}

Heterogeneous graph neural networks extend graph representation learning to graphs with multiple node and relation types. HAN \cite{wang2019han} uses node-level attention and semantic-level attention to aggregate information from meta-path-based neighborhoods. HGT \cite{hu2020hgt} introduces type-specific projections and attention mechanisms for heterogeneous graph transformers. MAGNN \cite{fu2020magnn} models intra-meta-path and inter-meta-path aggregation to capture fine-grained semantic dependencies. HeCo \cite{wang2021heco} further uses contrastive learning between network schema and meta-path views to improve heterogeneous graph representation learning.

These methods show the value of semantic aggregation, but they are usually designed for clean graphs. Under adversarial structural perturbations, their reliance on observed edges and meta-path-induced neighborhoods may cause corrupted semantics to be propagated. DVCL keeps the semantic modeling ability of heterogeneous graph learning while explicitly introducing topology purification and a feature-induced view for robustness.

\subsection{Robust Graph Neural Networks}

Robust graph neural networks aim to reduce the sensitivity of graph models to adversarial or noisy structures. GCN-Jaccard \cite{wu2019gcnjaccard} removes edges between nodes with low feature similarity. RGCN \cite{zhu2019rgcn} uses Gaussian distributions to model hidden representations and reduce adversarial effects. Pro-GNN \cite{jin2020prognn} jointly learns graph structure and GNN parameters with sparsity and low-rank constraints. GNNGuard \cite{zhang2020gnnguard} assigns edge weights according to feature similarity, and SimP-GCN \cite{jin2021simpgcn} incorporates similarity-preserving constraints into graph convolution.

Most robust GNNs are developed for homogeneous graphs. Directly applying them to heterogeneous graphs may ignore typed relations and semantic paths. DVCL differs by operating on meta-path-induced semantic structures and by using feature-induced evidence as an auxiliary view rather than as a simple edge-pruning rule.

\subsection{Robust Heterogeneous Graph Learning}

Robust heterogeneous graph learning has recently attracted attention because structural attacks can be amplified through heterogeneous semantics. Methods such as RoHe \cite{rohe}, HeteGuard \cite{heteguard}, and HSeCo \cite{hseco} attempt to defend HGNNs by modeling relation-aware perturbations, filtering unreliable semantic connections, or improving semantic topology learning.

These methods are closely related to DVCL, but they still mainly rely on the attacked heterogeneous topology when constructing candidate neighborhoods. DVCL addresses this limitation by adding a feature-induced graph that does not directly depend on perturbed edges. The cross-view objective then encourages topology-view representations and feature-view representations to agree, improving robustness when either view is noisy.

\subsection{Graph Contrastive Learning}

Graph contrastive learning learns representations by maximizing agreement between related graph views. DGI \cite{velickovic2019dgi} contrasts node representations with graph-level summaries. MVGRL \cite{hassani2020mvgrl} contrasts diffusion and adjacency views. GRACE \cite{zhu2020grace} and GCA \cite{zhu2021gca} generate graph augmentations and align node representations across augmented views. In heterogeneous graphs, HeCo \cite{wang2021heco} contrasts schema and meta-path views to learn semantic representations.

DVCL is also based on cross-view agreement, but its views are designed for robustness rather than only self-supervised representation learning. The topology view captures purified heterogeneous semantics, while the feature-induced view provides an attack-resistant reference. The contrastive objective is therefore used to improve consistency between semantic topology and feature similarity under structural attacks.

% ========== 6 Conclusion ==========
\section{Conclusion}

In this paper, we proposed DVCL, a dual-view contrastive learning framework for robust heterogeneous graph neural networks. DVCL constructs a purified topology view from meta-path-induced semantic structures and a feature-induced view from target-node similarity. The topology view preserves heterogeneous semantic information through meta-path construction, edge confidence estimation, semantic attention, and two-stage filtering, while the feature view provides complementary evidence that is less dependent on attacked graph structure. A symmetric cross-view contrastive loss aligns the two views and improves representation consistency under structural perturbations.

The first version of the experimental section defines the evaluation protocol for clean graphs, global structural attacks, and target attacks on heterogeneous graph benchmarks. After filling in the measured results, the final analysis will quantify the robustness gains of DVCL, the contribution of each component, and the sensitivity to key hyperparameters. Future work may extend DVCL to dynamic heterogeneous graphs and investigate more adaptive attacks that jointly manipulate topology and node attributes.

% ========== References ==========
\bibliographystyle{plain}
\bibliography{references}

\end{document}

